\documentclass[11pt]{article}
\usepackage[margin=1in]{geometry}
\usepackage{amsmath,amssymb,booktabs,longtable,array}
\usepackage[hidelinks]{hyperref}
\setlength{\parindent}{0pt}
\setlength{\parskip}{0.5em}
\newcommand{\R}{\mathbb{R}}
\newcommand{\ind}{\mathbf{1}}
\newcommand{\code}[1]{\texttt{\detokenize{#1}}}
\newcommand{\clip}{\operatorname{clip}}
\newcommand{\wrap}{\operatorname{wrap}}
\title{Mathematical Specification of the Auditorium Geometry\\and Its Interaction with Agents}
\author{}
\date{}
\begin{document}
\maketitle

\section{Scope and interpretation}

This specification describes the current two-plan Python implementation,
including the removal of familiarity from behavioral utility. It is a
specification of the implemented simulation, rather than a surveyed or
regulatory description of Bailey Hall. The seat inventories and connection
counts follow the supplied building description; metric dimensions, travel
lengths, visibility surfaces, and throughput coefficients are schematic.

The building is a hybrid spatial system: agents move continuously within
state-dependent planar regions and change regions or floors through discrete
transitions. It is not represented by one unrestricted polygon or by a fully
resolved three-dimensional mesh. In particular, overlap between two regions'
coordinate descriptions does not authorize an agent to change its region.

Coordinates and lengths below are in metres, time in seconds, velocities in
metres per second, and mathematical angles in radians. Numerical parameters
are a snapshot of \code{default_world.json} and
\code{default_policy_settings.json} at the time of writing. Later configuration
edits supersede the numerical values in this document.

\section{Coordinate system and independent floor plans}

Let the occupied floor labels be
\[
 \mathcal{L}=\{0,1\},
\]
where $0$ is the main plan containing orchestra, loge, corridor, and foyer,
and $1$ is the balcony plan. The elevated loge remains on plan $0$; the
orchestra-to-loge ascent is represented by explicit access connectors.
No continuous height coordinate or gravitational dynamics is assigned.

Both plans use the reference centre
\[
 \mathbf c=(0,7.5)^{\mathsf T}.
\]
For $\mathbf x=(x,y)^{\mathsf T}$ define
\begin{align}
 r(\mathbf x)&=\sqrt{x^2+(y-7.5)^2},\\
 \theta(\mathbf x)&=\operatorname{atan2}(y-7.5,x),\\
 \mathbf p(r,\theta)&=(r\cos\theta,\;7.5+r\sin\theta)^{\mathsf T}.
\end{align}
The principal auditorium construction occupies the semicircle
$0\leq\theta\leq\pi$. Increasing $y$ generally points toward the rear foyer
in the drawing; these coordinates are not georeferenced survey coordinates.
Angular fractions stored in the configuration are multiplied by $\pi$.

An agent's hybrid state may be written
\[
 S_i(t)=\bigl(\ell_i,q_i,z_i,\mathbf x_i,\mathbf v_i,a_i,
                  s_i,u_i,\mathcal M_i,\chi_i\bigr),
\]
where $\ell_i$ is its floor, $q_i$ its region, $z_i$ its seating zone,
$a_i$ its body radius, $s_i$ its initial seat, $u_i$ its current local target,
$\mathcal M_i$ its memory, and $\chi_i$ its transit or completion state.
The implementation uses floor $-1$ for transit and $-2$ for completed agents.
Identical planar coordinates on different floors do not imply interaction.

\section{Seating construction and population initialization}

\subsection{Section parameters}

\begin{center}
\small
\begin{tabular}{lrrrrrr}
\toprule
Section & Floor & Seats $C_s$ & Rows $J_s$ & First radius & Row spacing & Pitch\\
\midrule
Orchestra & 0 & 527 & 16 & 6.70 & 0.64 & 0.58\\
Loge      & 0 & 486 &  7 & 17.50 & 0.68 & 0.60\\
Balcony   & 1 & 316 &  9 & 16.00 & 0.68 & 0.67\\
\bottomrule
\end{tabular}
\end{center}

The total inventory is
\[
 C=527+486+316=1329.
\]
No additional wheelchair agents or wheelchair-specific motion classes are
introduced. Every occupied seat corresponds to one agent with the same
default body radius, although preferred speeds and behavioral weights vary.

For row index $j\in\{0,\ldots,J_s-1\}$,
\[
 R_{s,j}=R_{s,0}+j\Delta R_s.
\]
The six principal aisle bearings are
\[
 (\alpha_1,\ldots,\alpha_6)
 =\pi(0.025,0.215,0.405,0.595,0.785,0.975).
\]
The row-block construction uses the ordered boundary list
\[
 (\beta_0,\ldots,\beta_7)
 =\pi(0.02,0.025,0.215,0.405,0.595,0.785,0.975,0.98).
\]
The nominal aisle half-width is $h=1.3$, hence the nominal full width is $2.6$.

\subsection{Seat allocation within rows and blocks}

For adjacent bearings $(\beta_b,\beta_{b+1})$, define the usable row length
and provisional seat capacity by
\begin{align}
 L_{s,j,b}&=R_{s,j}(\beta_{b+1}-\beta_b)-2h,\\
 M_{s,j,b}&=\max\!\left\{0,
             \left\lfloor\frac{L_{s,j,b}}{\text{pitch}_s}\right\rfloor+1\right\}.
\end{align}
Blocks with $M_{s,j,b}=0$ are omitted. Let
$M_s=\sum_{j,b}M_{s,j,b}$. The actual allocation begins with
\[
 n_{s,j,b}=\left\lfloor C_s\frac{M_{s,j,b}}{M_s}\right\rfloor.
\]
Remaining seats are assigned to the largest fractional remainders, breaking
ties by row and block index. Consequently $\sum_{j,b}n_{s,j,b}=C_s$ exactly.
This requires $M_s\geq C_s$.

Seat-centre bearing limits within an occupied block are
\[
 \theta^-_{s,j,b}=\beta_b+\frac{h}{R_{s,j}},\qquad
 \theta^+_{s,j,b}=\beta_{b+1}-\frac{h}{R_{s,j}}.
\]
For $n=n_{s,j,b}>1$ and $k=0,\ldots,n-1$,
\[
 \theta_{s,j,b,k}=\theta^-_{s,j,b}
     +\frac{k}{n-1}(\theta^+_{s,j,b}-\theta^-_{s,j,b}),\qquad
 \mathbf x_{s,j,b,k}=\mathbf p(R_{s,j},\theta_{s,j,b,k}).
\]
A single seat is placed at the angular midpoint. Pitch controls provisional
capacity; final centres are evenly spread over each allocated block and do
not necessarily have exactly that pitch.

\subsection{Occupancy for a requested population}

For integer $1\leq n\leq1329$, allocate agents proportionally:
\[
 n_s^{(0)}=\left\lfloor n\frac{C_s}{1329}\right\rfloor.
\]
Largest fractional remainders distribute the remaining agents among sections.
Seats are sampled without replacement within each section using the seeded
generator. At $n=300$, the counts are 119 orchestra, 110 loge, and 71 balcony.
Agent identity and initial seat identity persist across all transfers.

\section{State-dependent navigable regions}

\subsection{Seating-zone aisle extents}

\begin{center}
\begin{tabular}{llrrl}
\toprule
Zone & Floor & $R_z^-$ & $R_z^+$ & Outward continuation\\
\midrule
Orchestra & 0 & 6.30 & 16.525 & Access connector to loge\\
Loge      & 0 & 17.30 & 22.20 & Main circulation rim\\
Balcony   & 1 & 15.70 & 22.75 & Balcony circulation rim\\
\bottomrule
\end{tabular}
\end{center}

For an agent assigned aisle bearing $\alpha_i$, the aisle centre-domain is
\[
 \mathcal I_i=\left\{\mathbf x:
 |r(\mathbf x)(\theta(\mathbf x)-\alpha_i)|
       \leq h-a_i-0.02,\quad
 R_{z_i}^- -0.08\leq r(\mathbf x)\leq R_{z_i}^++0.03\right\}.
\]
The lateral expression is an arc-based approximation, rather than the exact
perpendicular distance to a radial line.

\subsection{Seat-row motion domains}

For an agent whose initial row radius is $R_i$ and row-end bearings are
$\theta_i^-,\theta_i^+$, the row acceptance predicate is
\[
 |r(\mathbf x)-R_i|<0.07,\qquad
 \theta_i^- -0.012\leq\theta(\mathbf x)\leq\theta_i^++0.012,
\]
on the initial seat's floor. The angular tolerance $0.012$ is in radians.
Actual accepted row motion is projected onto $r=R_i$ as described below.
Unoccupied seats are not individual solid obstacles; row restrictions prevent
radial shortcuts through seating blocks.

\subsection{Curved circulation regions}

Write
\[
 \mathcal A(R^-,R^+,\gamma^-,\gamma^+)
 =\{\mathbf x:R^-\leq r(\mathbf x)\leq R^+,\;
                     \pi\gamma^-\leq\theta(\mathbf x)\leq\pi\gamma^+\}.
\]
The configured regions are
\begin{center}
\begin{tabular}{lrrrrr}
\toprule
Region & Floor & $R^-$ & $R^+$ & $\gamma^-$ & $\gamma^+$\\
\midrule
Main rim    & 0 & 22.02 & 24.05 & 0.012 & 0.988\\
Corridor    & 0 & 24.00 & 27.30 & 0.005 & 0.995\\
Balcony rim & 1 & 22.57 & 24.50 & 0.008 & 0.992\\
\bottomrule
\end{tabular}
\end{center}
These annular bounds constrain body centres directly. Unlike the rectangular
regions below, the implementation does not uniformly erode these annuli by
body radius. Additional clearance enters waypoint placement, obstacle forces,
and doorway checks. Thus this is not a uniform disc-eroded floor polygon.

\subsection{Rectangular regions}

\begin{center}
\begin{tabular}{lrrrrr}
\toprule
Region & Floor & $x^-$ & $x^+$ & $y^-$ & $y^+$\\
\midrule
Foyer              & 0 & -12.40 & 12.40 & 32.20 & 41.20\\
East stage landing & 0 & 21.50  & 24.25 & 3.80 & 6.80\\
West stage landing & 0 & -24.25 & -21.50 & 3.80 & 6.80\\
\bottomrule
\end{tabular}
\end{center}
Away from an exit apron, an agent centre must satisfy
\[
 x^-+a_i\leq x_i\leq x^+-a_i,\qquad
 y^-+a_i\leq y_i\leq y^+-a_i.
\]
Floor and region labels must also match. The three rectangular patches and
the annular regions are linked by the transition graph, not by unrestricted
movement through their geometric union.

\section{Orchestra--loge barrier}

The nominal barrier has centre radius $R_B=16.9$, radial thickness $b_B=0.35$,
and six openings of width $w_B=1.3$ at bearings $\alpha_1,\ldots,\alpha_6$.
Its occupied radial band is $[16.725,17.075]$ within the auditorium semicircle.

For a main-floor agent, the implementation declares radial body overlap when
\[
 |r_i-R_B|<\frac{b_B}{2}+a_i.
\]
Such a position is accepted only when an opening fits the body:
\[
 \exists k\in\{1,\ldots,6\}:\quad
       r_i|\theta_i-\alpha_k|\leq\frac{w_B}{2}-a_i-0.01.
\]
This barrier predicate supplements the region predicate. It does not create
a continuous orchestra-to-loge walking surface: a successful access admission
initiates the timed transfer defined later.

\section{Portal topology and coordinates}

\subsection{Directed hybrid graph}

Let $\mathcal G=(\mathcal V,\mathcal P)$ be a directed multigraph of region
states and portals. A portal contains its source floor, source region,
optional source zone, width, position, type, availability flag, and either a
named destination entrance or an exterior exit identity.

The principal paths are
\begin{align*}
 \text{orchestra row}&\longrightarrow\text{orchestra aisle}
 \xrightarrow{\mathrm{ACCESS}}\text{loge aisle},\\
 \text{loge row}&\longrightarrow\text{loge aisle}
 ,\\
 \text{loge aisle}&\longrightarrow\text{main rim}
 \xrightarrow{\mathrm{AUD}}\text{corridor},\\
 \text{corridor}&\xrightarrow{\mathrm F}\text{foyer}
 \xrightarrow{\mathrm M}\text{outside},\\
 \text{corridor}&\xrightarrow{\mathrm{SIDE}}\text{outside},\\
 \text{balcony row}&\longrightarrow\text{balcony aisle}
 \longrightarrow\text{balcony rim}
 \xrightarrow{\mathrm{BACK}}\text{foyer},\\
 \text{balcony rim}&\xrightarrow{\mathrm{STAGE}}
       \text{first-floor stage landing}
 \xrightarrow{\mathrm{GROUND\_STAGE}}\text{outside}.
\end{align*}

There are 34 configured portals: six access connectors, nine auditorium doors,
three foyer-entry doors, four balcony stair connectors, eleven exterior
thresholds, and one excluded elevator. There are 22 named arrival entrances.
Exterior doors do not require a destination entrance.

\subsection{Parametric portal locations}

\begin{longtable}{p{0.20\linewidth}p{0.56\linewidth}r}
\toprule
Family & Position or bearing definition & Width\\
\midrule
\endhead
ACCESS1--6 & $\mathbf p(16.525,\alpha_k)$ & 1.30\\
AUD0--8 & $\mathbf p(23.25,\pi(0.06+0.11j))$, $j=0,\ldots,8$ & 1.35\\
F0--2 & $\mathbf p(25.6,\pi(0.44+0.06j))$, $j=0,1,2$ & 1.65\\
M0--4 & $((j-2)2.35,41.2)^{\mathsf T}$, $j=0,\ldots,4$ & 1.75\\
SIDE E1, E2, W2, W1 & $\mathbf p(27.3,\pi f)$ for $f\in\{0.15,0.34,0.66,0.85\}$ & 1.35\\
BACK E, W & $\mathbf p(23.15,\pi f)$ for $f\in\{0.35,0.65\}$ & 1.60\\
STAGE E, W & $\mathbf p(23.15,\pi f)$ for $f\in\{0.025,0.975\}$ & 1.15\\
GROUND STAGE E, W & $(24.25,5.3)^{\mathsf T}$ and $(-24.25,5.3)^{\mathsf T}$ & 1.15\\
ADA W & $(-11.2,39.0)^{\mathsf T}$; excluded from evacuation & 1.20\\
\bottomrule
\end{longtable}

The five M thresholds share the reported exit category \code{MAIN}. The
balcony STAGE portals are transfers, whereas GROUND STAGE portals are final
thresholds. A shared identifier in the entrance and portal namespaces does
not imply shared coordinates or immediate completion.

\subsection{Named destination entrances}

The access destinations are $\mathbf p(17.42,\alpha_k)$ in loge aisles.
AUD destinations are $\mathbf p(24.66,\phi_j)$ in the corridor, where
$\phi_j=\pi(0.06+0.11j)$. Both families have outward radial arrival normals.
Foyer-entry destinations retain their corresponding F-door $x$ coordinate
and have $y=32.85$, with arrival normal $(0,1)^{\mathsf T}$.

BACK E and W arrive at $(10.6,35.3)^{\mathsf T}$ and
$(-10.6,35.3)^{\mathsf T}$, respectively, with normal $(0,1)^{\mathsf T}$.
STAGE E and W arrive at $(23.1,5.3)^{\mathsf T}$ and
$(-23.1,5.3)^{\mathsf T}$, with normals $(1,0)^{\mathsf T}$ and
$(-1,0)^{\mathsf T}$. These arrivals are still inside the first-floor plan.

\section{Agent bodies and geometric acceptance}

The default radius is fixed:
\[
 a_i=0.2025\qquad(i=1,\ldots,n).
\]
The model treats each active planar body as a disc for interpersonal exclusion.
For two agents on the same floor, a candidate centre $\mathbf x$ must satisfy
\[
 \|\mathbf x-\mathbf x_j\|^2\geq1.005(a_i+a_j)^2.
\]
With equal default radii this is approximately $0.406$ metres of centre
separation. Agents on other floors, agents in transit, and completed agents
are absent from this planar collision test.

Let $C_i(\mathbf x;q_i,z_i,\ell_i)$ denote the implemented state-dependent
region/barrier/apron predicate. The admissible endpoint set at a sequential
update is
\[
 \mathcal F_i(t)=\left\{\mathbf x:C_i(\mathbf x;q_i,z_i,\ell_i)=1,
 \ \|\mathbf x-\mathbf x_j\|^2\geq1.005(a_i+a_j)^2
 \text{ for active }j\ne i\text{ on floor }\ell_i\right\}.
\]
Exit enablement and capacity add crossing constraints. Region enforcement,
hard-body exclusion, and door capacity are enabled in the current world.
Ordinary body/wall checks are primarily endpoint checks, not complete swept
disc collision detection; exterior threshold intersections are swept tests.

\section{Local actions, visibility, and route choice}

\subsection{Available and perceived targets}

An action identifies a next local target: a row end, aisle continuation,
eligible door or connector, or a short exploration waypoint. Portal
eligibility requires source floor, source region, evacuation permission,
and any required source zone to match. An access connector also requires
$|\phi_p-\alpha_i|<10^{-6}$.

Newly observed targets must satisfy
\[
 \|\mathbf p_a-\mathbf x_i\|\leq R_{\rm view},\qquad
 \widehat{\mathbf h}_i\cdot
 \frac{\mathbf p_a-\mathbf x_i}{\|\mathbf p_a-\mathbf x_i\|}
 \geq\cos(\Phi/2),
\]
and the configured line of sight must be unblocked. Current values are
$R_{\rm view}=12$ and $\Phi=2\pi$, permitting an all-around scan.
Nearby-person observation has the separate range $R_{\rm obs}=4$.

Visibility uses the barrier surfaces and openings, a schematic auditorium
wall at corridor inner radius $24.0$ with enabled AUD apertures, foyer walls
with F/M apertures, and optional floor-specific line segments in
\code{visibilityWalls}. Body and seat occlusion and eye height are not modeled.
Additional visibility segments do not create physical collision walls.

Previously observed targets, the current target, and locations remembered
from admission may remain available outside instantaneous sight. Recorded
admission paths can be supplied by seat; otherwise a path from M2 is
reconstructed through reverse portal topology. This initialization does not
simulate a pre-emergency admission crowd. Remembering a target gives no
familiarity reward in utility.

When no door is seen or remembered in an annular region, short local left/right
search targets are generated from region bounds, without using hidden door
locations. With no action elsewhere, the stationary scan heading rotates on
the next decision. Known closed portals are excluded. Closure discovery occurs
within $2.2$ metres when the eligible portal is visible.

\subsection{Distance and behavioral features}

Distance to a row end is the remaining row arc length. For an annular portal,
the approximation is
\[
 d_i(a,t)=r_i|\wrap(\phi_a-\theta_i)|+|r_i-r(\mathbf p_a)|,
\]
where $\wrap(\psi)=\operatorname{atan2}(\sin\psi,\cos\psi)$.
An ordinary aisle continuation uses $\max(0,R_z^+-r_i)$; other actions use
Euclidean target distance. These are local travel estimates, not globally
optimal path lengths.

The current four-dimensional behavioral feature vector is
\[
 \mathbf X_i(a,t)=(-T_i(a,t),\;P_i(a,t),\;-N_i(a,t),\;-C_i(a,t))^{\mathsf T}.
\]
There is no familiarity feature or weight. Perceived time is
\[
 T_i(a,t)=\frac{1}{\tau_T}\left[
 \frac{d_i(a,t)}{\max(0.35,v_i^0)}+Q_i(a,t)+B_i(a,t)\right].
\]
The current normalization is $\tau_T=1$ second. When enabled, $Q_i$ estimates
delay from observed earlier-priority people waiting ahead at that portal,
using their count divided by $\max(0.1,1.3w_p)$. Waiting means speed at most
$0.25$ within $3$ metres of the selected portal. Priority follows first-wait
time, with agent ID breaking exact ties.

$B_i$ is an optional opposing-row blockage estimate. If the agent has been
physically blocked since $t_b$ and observes an opposing person ahead in the
same seat row and block, the obstructed action receives
\[
 B_i(a,t)=\kappa_B\bigl(\max(0,t-t_b)+\tau_{\rm dec}\bigr),\qquad\kappa_B=1.
\]
Otherwise it is zero. This extrapolation from experienced waiting is a
heuristic, not a measured future delay. Disabling perceived queue delay
disables both $Q_i$ and $B_i$.

$P_i$ is the fraction of observed nearby agents whose intended local direction
is within $45$ degrees of the candidate direction; stationary agents with
a target can count. $N_i$ is the maximum of an opposing-motion fraction and
a queue-cutting fraction. The latter counts earlier-priority waiters whom
the candidate would place behind the deciding agent, using a clearance
margin $a_i+a_j+0.05$ in target distance. It therefore includes stationary
queues. The opposing-motion component considers speeds above $0.12$ and
directional cosine below $-0.15$. These are operational norm definitions.

The switching feature is
\[
 C_i(a,t)=\ind\{\text{a current target exists in this stage and }a\ne u_i\}.
\]
Beginning a new physical stage resets this commitment. Feature observations
are frozen at tick start; the same feature-construction rules apply to all
agents, independently of their utility weights.

\subsection{Weights and decision rule}

Each agent receives fixed weights in the order $(T,P,N,S)$:
\begin{align}
 \mathbf z_i&\sim\mathcal N_4(\boldsymbol\mu,\Sigma),\qquad
 w_{ij}=e^{z_{ij}},\\
 \Sigma&=\sigma^2[(1-\rho)I_4+\rho\mathbf1_4\mathbf1_4^{\mathsf T}],
 \qquad\boldsymbol\mu=-\tfrac12\sigma^2\mathbf1_4.
\end{align}
Thus $\mathbb E[\mathbf w_i]=\mathbf1_4$, with strictly positive, unbounded,
unnormalized components. Positive definiteness requires $-1/3<\rho<1$.
Current values are $\sigma=1$ and $\rho=-0.15$.

Behavioral utility is
\[
 U_i^{\rm beh}(a,t)=-w_{iT}T_i+w_{iP}P_i-w_{iN}N_i-w_{iS}C_i.
\]
The default full utility additionally retains the earlier distance and noise
terms:
\[
 U_i(a,t)=\left(\frac{d_{i,\min}(t)+\ell_0}{d_i(a,t)+\ell_0}\right)^{\beta_i}
       +U_i^{\rm beh}(a,t)
       +\pi_i(1-e^{-t/60})\delta_{ia}(t)+\epsilon_{ia}(t),
 \qquad \ell_0=1\text{ metre}.
\]
Here $\beta_i\in[0,1]$ is the bounded distance-sensitivity trait,
$\pi_i\in[0,1]$ is panic sensitivity, and the two shock types use independent
zero-mean normal draws of standard deviation $0.2$, held within 15-second
bins. There is no standalone global crowd reward. The
\code{behavioralOnly} option instead uses $U_i^{\rm beh}$ alone.

After its reaction delay, an agent chooses
\[
 u_i(t)\in\operatorname*{arg\,max}_{a\in\mathcal A_i(t)} U_i(a,t),
\]
with first-enumerated action breaking exact ties. The current reconsideration
interval is $\tau_{\rm dec}=3$ seconds in every decision region; physical stage
changes and discovered target closures can cause earlier reconsideration.
Consequently the nearest door need not maximize utility.

\section{Geometry-dependent waypoints and social forces}

\subsection{Local guidance}

The selected action determines a waypoint $\mathbf g_i$, rather than permission
to move along a straight line through walls. Row waypoints are row ends;
aisle waypoints follow the selected aisle outward. Near a narrow orchestra
opening, a sideways approach is aligned at $R_z^+-0.35$ before advancing.

In an annulus, define
\[
 b_{ip}=\max(0.02,w_p/2-a_i-0.08),\qquad
 \Delta\theta=\wrap(\phi_p-\theta_i).
\]
Direct approach requires $|\Delta\theta|<\pi/2$ and
$|r_i\sin\Delta\theta|\leq b_{ip}$. Otherwise the safe circulation radius is
\[
 r_i^{\rm safe}=\clip(r_i,R^-+a_i+0.06,R^+-a_i-0.06).
\]
A radial correction is used when $|r_i-r_i^{\rm safe}|>0.05$; otherwise the
waypoint advances angularly by $\clip(\Delta\theta,-0.065,0.065)$.
This allows different circulation radii for opposing flows.

An aligned exterior waypoint extends $0.6$ metres beyond its threshold.
For a rectangular exit approached from the side, guidance first retreats
inside the wall to depth $a_i+0.08$, if needed, and then aligns with the chosen
opening. This permits a late target change from another doorway's apron.

\subsection{Continuous acceleration proposal}

Preferred speed $v_i^0$ is sampled using the configured normal speed model
(mean $1.29$, standard deviation $0.19$, bounds $0.5$ and $2.3$).
Reaction delays use the configured lognormal model with arithmetic mean $2$,
standard deviation $1.2$, and bounds $0$ and $40$. These distribution bounds
are clipping limits; the clipped sample need not retain the unbounded
distribution's exact moments.
Let $\widehat{\mathbf e}_i$ point from the current position to its waypoint.
The manuscript-force proposal is
\begin{align}
 \frac{d\mathbf v_i}{dt}
   &=\mathbf f_i^{\rm des}+\mathbf f_i^{\rm obs}+\mathbf f_i^{\rm int},\\
 \mathbf f_i^{\rm des}
   &=\frac{v_i^0\widehat{\mathbf e}_i-\mathbf v_i}{\tau_r},
      \qquad\tau_r=0.65,\\
 \mathbf f_i^{\rm obs}
   &=\sum_{w\in W_i}A\exp(-d_{iw}/r_i^c)\widehat{\mathbf n}_{iw},
      \qquad A=0.8,\quad r_i^c=0.15\text{ by default}.
\end{align}
At the exact waypoint, desired velocity is zero and relaxation brakes the
agent. Individual comfort distances can override $r_i^c$. Obstacle distances
are nonnegative centre-to-represented-surface distances. $W_i$ contains the
schematic region or aisle boundary contacts and the orchestra barrier, rather
than every real seat or wall. Open exterior apertures suppress the artificial
wall contribution across their opening.

For same-floor neighbors within $2.6$ metres, set
\begin{align}
 d_{ij}&=\|\mathbf x_j-\mathbf x_i\|,\qquad
 \widehat{\mathbf x}_{ij}=(\mathbf x_j-\mathbf x_i)/d_{ij},\\
 \mathbf D_{ij}&=\lambda(\mathbf v_i-\mathbf v_j)+\widehat{\mathbf x}_{ij},
 \qquad\widehat{\mathbf D}_{ij}=\mathbf D_{ij}/\|\mathbf D_{ij}\|,\\
 \widehat{\mathbf N}_{ij}&=(-\widehat D_{ij,y},\widehat D_{ij,x})^{\mathsf T},\\
 \vartheta_{ij}&=\operatorname{atan2}
   (\widehat{\mathbf D}_{ij}\mathbin{\times}\widehat{\mathbf x}_{ij},
    \widehat{\mathbf D}_{ij}\cdot\widehat{\mathbf x}_{ij}).
\end{align}
The two-dimensional cross product here is scalar. Interpersonal acceleration is
\[
 \mathbf f_i^{\rm int}=-\sum_j e^{-d_{ij}/\|\mathbf D_{ij}\|}
 \left[
 e^{-(n'\vartheta_{ij}\|\mathbf D_{ij}\|)^2}\widehat{\mathbf D}_{ij}
 +\operatorname{sgn}(\vartheta_{ij})
  e^{-(n\vartheta_{ij}\|\mathbf D_{ij}\|)^2}\widehat{\mathbf N}_{ij}
 \right],
\]
with $\lambda=0.25$ and $n=n'=1$. Coincident positions and vanishing
$\mathbf D_{ij}$ are skipped at tolerance $10^{-12}$ to avoid undefined
directions. These force calculations have no familiarity multiplier.

\subsection{Discrete movement and local avoidance}

For $\Delta t=0.1$, explicit Euler first gives
\[
 \widetilde{\mathbf v}_i^{k+1}
 =\mathbf v_i^k+\Delta t\,
       (\mathbf f_i^{\rm des}+\mathbf f_i^{\rm obs}+\mathbf f_i^{\rm int}).
\]
Its magnitude is capped by $1.35v_i^0$ and, for a nonzero waypoint distance,
by $\|\mathbf g_i-\mathbf x_i\|/\Delta t$.

Optional local avoidance checks predicted closest approach over $0.65$ seconds.
It uses up to $0.04$ metres of extra padding, reduced at existing close contact.
When the original request is unsafe or physical progress stalls, it considers
zero velocity and nearby directions at full and half preferred speed. The
angles are $0$, $\pm20$, $\pm40$, $\pm65$, $\pm90$, $\pm120$, and $180$ degrees.
A deterministic passing-side preference resolves head-on symmetry. Candidate
scores reward forward progress and passing, with a small penalty for departure
from the force proposal. Progress of $0.12$ metres resets the stall clock;
two seconds of insufficient progress activates the stalled condition.
This stage does not change the selected target and bypasses seat-row motion.

Let $\mathbf v_i^{\rm req}$ be the resulting request. Its displacement is
\[
 \Delta\mathbf x_i=\mathbf v_i^{\rm req}\Delta t,
\]
rescaled if necessary to have magnitude at most $0.7$ metres per tick.
For a row agent, define $\mathbf e_\theta=(-\sin\theta_i,\cos\theta_i)^{\mathsf T}$;
the candidate position is projected by
\[
 \theta_i^{\rm cand}=\clip\left(
 \theta_i+\frac{\mathbf e_\theta\cdot\Delta\mathbf x_i}{R_i},
 \theta_i^-,\theta_i^+\right),\qquad
 \mathbf x_i^{\rm cand}=\mathbf p(R_i,\theta_i^{\rm cand}).
\]
Other agents first try the full displacement. If it fails, the engine tries,
in order, half displacement, its $x$ component alone, and its $y$ component
alone. A valid candidate is accepted; otherwise position is retained and
realized velocity is set to zero. Accepted realized velocity is displacement
divided by $\Delta t$. These constraints and fallbacks mean actual motion is
not simply an unconstrained solution of the force differential equation.

\section{Ordinary region changes and doorway service}

Near a row end, within $0.07$ metres of row arc length, the policy requests an
aisle landing at bearing
\[
 \alpha_i+s_i\frac{h-a_i-0.06}{R_i},\qquad s_i\in\{-1,+1\}.
\]
This landing must pass geometry and body-clearance checks. An aisle agent
requests its outward continuation when $r_i\geq R_z^+-0.16$. Orchestra agents
must use an ACCESS connector; loge and balcony agents change to their respective
rims. Ordinary region changes cannot alter the floor or seating zone and are
limited to the allowed row-to-aisle and aisle-to-rim transitions.

For an interior portal $p$, admission requires eligibility, enablement,
\[
 \|\mathbf x_i-\mathbf p_p\|\leq0.65,
\]
and, for a radial portal, tangential clearance
\[
 |(-\sin\phi_p,\cos\phi_p)\cdot(\mathbf x_i-\mathbf p_p)|
       \leq w_p/2-a_i-0.01.
\]
Door capacity additionally requires an available service token.

Let $b_p$ be the token count. Initially and at each enabled tick,
\begin{align}
 b_p(0)&=\max(1,1.2w_p),\\
 b_p^{\rm pre}&=\min\{\max(1,1.35w_p),\ b_p+1.3w_p\Delta t\}.
\end{align}
Each successful admission or exterior crossing consumes one token. These
empirical width-dependent coefficients define the implemented service model;
they are not measured or code-certified discharge rates. Disabled portals
do not refill.

An interior ordinary door performs a discrete transfer to a valid named
destination entrance during the same tick. It does not integrate a continuous
doorway passage. If every candidate landing is occupied, no transfer occurs
and no token is consumed. Arrival updates the destination floor, region, zone,
and position, preserves identity and weights, and sets planar velocity to zero.

\section{Timed access and staircase connectors}

\begin{center}
\begin{tabular}{lrrrr}
\toprule
Connector & Total length & Flights & Lanes & Divider width\\
\midrule
Each ACCESS & 1.5 & $(1.5)$ & 1 & 0\\
Each BACK stair & 15 & $(7.5,7.5)$ & 2 & 0.1\\
Each STAGE stair & 12 & $(12)$ & 1 & 0\\
\bottomrule
\end{tabular}
\end{center}

For $m_p$ lanes and divider width $d_p^{\rm div}$, the lane width is
\[
 w_p^{\rm lane}=\frac{w_p-d_p^{\rm div}}{m_p}.
\]
Each BACK lane is therefore $0.75$ metres wide. Agent $i$ can enter only if
$w_p^{\rm lane}>2a_i+0.02$. Connector speed, entry headway, and the lane
occupancy bound used for admission are
\begin{align}
 v_{ip}^{\rm tr}&=\max(0.4,\eta_p v_i^0),\qquad
 \eta_p=\begin{cases}0.7,&p\text{ is a stair},\\1,&p\text{ is ACCESS},\end{cases}\\
 h_{ip}&=\max\!\left(0.35,\frac{2a_i+0.1}{v_{ip}^{\rm tr}}\right),\\
 K_{ip}&=\max\!\left(1,\left\lfloor\frac{L_p}{2a_i+0.1}\right\rfloor\right).
\end{align}
A lane is eligible if its next-entry time has passed and its current occupancy
is below $K_{ip}$. Among eligible lanes the model chooses minimum occupancy,
then earliest next-entry time, then lowest lane index.

For admission at $t_e$, nominal flight completion times are
\[
 t_{i,f}^{\rm nominal}=t_e+
       \sum_{g=1}^{f}\frac{L_{p,g}}{v_{ip}^{\rm tr}}.
\]
The final scheduled arrival is
\[
 t_i^{\rm sched}=\max\left(t_{i,F}^{\rm nominal},
                        t_{\rm previous}^{\rm sched}+h_{ip}\right).
\]
The lane's next-entry time becomes $t_e+h_{ip}$. Each lane is FIFO. An agent
can arrive only after its scheduled time and when it is first in its lane.
A blocked destination keeps it in transit, blocking later arrivals in that
lane until a valid landing is available. Disabling a source portal prevents
new admissions; it does not cancel agents already traveling in that connector.

For destination centre $\mathbf e_p$, arrival normal $\mathbf n_p^{\rm in}$,
and tangent $\mathbf t_p^{\rm in}=(-n_y^{\rm in},n_x^{\rm in})$, lane
$\ell\in\{0,\ldots,m_p-1\}$ tries
\[
 \mathbf x^{\rm land}_{i,\ell,d}
 =\mathbf e_p+\left(\ell-\frac{m_p-1}{2}\right)w_p^{\rm lane}\mathbf t_p^{\rm in}
                +d\mathbf n_p^{\rm in},\qquad
 d\in\{0,0.25,0.5,0.75\}.
\]
The first geometrically valid and unoccupied landing is selected. Ordinary
doors use the same landing construction with one lane and no lateral offset.

During timed transit, the agent is removed from both planar body indexes.
Its identity and preferences persist, but no continuous stair position,
stair-body contact force, or interaction between separate lanes is computed.
Leaving the balcony is therefore a change of plane, not completion.

\section{Exterior thresholds and exact completion}

For an exterior portal $p$, let $\mathbf p_p$ be its centre,
$\mathbf n_p$ its outward unit normal, and
$\mathbf t_p=(-n_{p,y},n_{p,x})^{\mathsf T}$. The geometric doorway segment is
\[
 \Gamma_p=\{\mathbf p_p+s\mathbf t_p:|s|\leq w_p/2\}.
\]
Normals are radial outward for side doors, $(0,1)^{\mathsf T}$ for M doors,
and $(\pm1,0)^{\mathsf T}$ for the corresponding ground stage thresholds.

For old and proposed centres $\mathbf x^0,\mathbf x^1$, set
\[
 s_0=\mathbf n_p\cdot(\mathbf x^0-\mathbf p_p),\qquad
 s_1=\mathbf n_p\cdot(\mathbf x^1-\mathbf p_p).
\]
An outward crossing requires $s_0<0\leq s_1$, with the implementation also
handling a start on the threshold using a $10^{-9}$ numerical tolerance.
The crossing fraction and point are
\[
 \lambda_* = \clip\left(\frac{-s_0}{s_1-s_0},0,1\right),\qquad
 \mathbf x_*=\mathbf x^0+\lambda_*(\mathbf x^1-\mathbf x^0).
\]
Body fit additionally requires
\[
 |\mathbf t_p\cdot(\mathbf x_*-\mathbf p_p)|\leq w_p/2-a_i-0.01.
\]
The portal must be eligible, enabled, have a token when capacity is enforced,
and be on floor $0$. A short exit apron permits body centres near and beyond
the threshold:
\[
 -a_i-0.05\leq\mathbf n_p\cdot(\mathbf x-\mathbf p_p)\leq0.8,
 \qquad |\mathbf t_p\cdot(\mathbf x-\mathbf p_p)|\leq w_p/2-a_i-0.01.
\]
An exterior candidate beyond the nominal boundary must contain a valid swept
threshold crossing; an endpoint in an apron does not authorize a diagonal
shortcut through the adjacent wall.

For a step from $t_k$ to $t_{k+1}=t_k+\Delta t$, the recorded exit time is
\[
 \tau_i=t_k+\lambda_*\Delta t.
\]
The centre is frozen at $\mathbf x_*$, the velocity is set to zero, and the
agent becomes inactive on the outside state. Proximity alone never counts as
completion. If several thresholds intersect a step, the earliest crossing is
used. Outdoor monumental and portico stairs lie after this completion event
and do not contribute additional modeled evacuation time. The ADA elevator
is permanently excluded from evacuation eligibility.

\section{Update order, conservation, and model limits}

Each tick advances the clock, replenishes enabled portal tokens, freezes
behavioral observations, and processes agents in a seeded random order.
An active planar agent observes and decides when due, requests motion,
undergoes geometric acceptance, and may request a region or portal transition.
Transit agents instead advance their lane and arrival state. Later physical
updates see earlier accepted positions in the same tick, even though behavioral
features use the frozen observation snapshot.

For all times,
\[
 N_0(t)+N_1(t)+N_{\rm transit}(t)+N_{\rm exited}(t)=n.
\]
Transfers preserve agents rather than creating replacement agents. The run
stops at full clearance or the configured time limit, currently 1800 seconds.
The reported clearance time is
\[
 T_{\rm clear}=\max_i\tau_i
\]
only if every agent exits; otherwise it is undefined and exported as null.

The main limitations are schematic metric geometry; state-dependent and
nonuniform wall clearance; endpoint-based ordinary collision checks; discrete
door transfers; timed rather than spatial stair motion; approximate sight and
norm features; and uncalibrated force, queue, and service parameters. The
model does not guarantee that every heterogeneous crowded realization clears.
No exact as-built reconstruction or evacuation certification is implied.

\section{Implementation correspondence}

\begin{longtable}{p{0.33\linewidth}p{0.60\linewidth}}
\toprule
Source & Mathematical responsibility\\
\midrule
\endhead
\code{default_world.json} & Dimensions, seat inventory, regions, portal records, named entrances, and world limits.\\
\code{geometry.py} & Polar transforms, largest-remainder seat construction, geometry dispatch.\\
\code{plans.py} & Region/barrier/apron predicates, portal eligibility, exterior intersections, local waypoints, obstacle contacts.\\
\code{engine.py} & Population allocation, collision checks, discrete motion, token refill, region changes, and completion.\\
\code{transfers.py} & Interior doorway landings, connector admission, lane scheduling, FIFO arrival, and threshold finalization.\\
\code{perception.py} & Sight range and angle, schematic wall/aperture occlusion, optional visibility segments.\\
\code{arrival.py} & Recorded or reconstructed admission-route knowledge.\\
\code{behavior.py} & Perceived local actions, frozen observations, four behavioral features, choice, and reconsideration.\\
\code{social_force.py}, \code{navigation.py} & Force proposal and optional collision-aware local steering.\\
\code{default_policy_settings.json} & Numerical behavioral, sensing, force, and steering coefficients.\\
\code{plan_viewer.py}, \code{viewer.py} & Read-only visualization; drawing primitives do not define collision acceptance.\\
\bottomrule
\end{longtable}

\end{document}
